% --- Archon physics formalization source begin ---
% archon:physics
% archon:covers PhyXMiniProblems/problem_phyx_mini_0240.lean
% archon:source-report reports/phyx_mini/problem_phyx_mini_0240.source.json

\chapter{Physics problem phyx\_mini\_0240}
\label{ch:PhyXMiniProblems_problem_phyx_mini_0240}

\paragraph{Problem source.}
\#\# Physical scenario

The string shown in figure is driven at a frequency of \textbackslash{}( 5.00\textbackslash{},\textbackslash{}mathrm\{Hz\} \textbackslash{}). The amplitude of the motion is \textbackslash{}( A = 12.0\textbackslash{},\textbackslash{}mathrm\{cm\} \textbackslash{}), and the wave speed is \textbackslash{}( v = 20.0\textbackslash{},\textbackslash{}mathrm\{m/s\} \textbackslash{}). Furthermore, the wave is such that \textbackslash{}( y = 0 \textbackslash{}) at \textbackslash{}( x = 0 \textbackslash{}) and \textbackslash{}( t = 0 \textbackslash{}).

\#\# Question

Determine the maximum transverse acceleration of an element of the string.

\#\# Answer choices

- A: A: \$128\textbackslash{}

- B: B: \$113\textbackslash{}

- C: C: \$108\textbackslash{}

- D: D: \$118\textbackslash{}

\#\# Figure caption (auxiliary; use the image as primary evidence)

The image depicts a transverse wave on a rope. 

- The rope originates from the left where it is attached to a wall or some fixed point.
- The rope then forms a sinusoidal wave moving to the right.
- An arrow labeled with \textbackslash{}(\textbackslash{}vec\{v\}\textbackslash{}) indicates the direction of the wave's propagation, pointing right.
- The wave's amplitude is marked with the label \textbackslash{}(A\textbackslash{}), shown as a vertical arrow from the equilibrium position (dashed horizontal line) to the peak of the wave.
- The wave pattern consists of alternating crests and troughs, characteristic of transverse waves.

\#\# Dataset metadata

Wave Properties; Physical Model Grounding Reasoning, Numerical Reasoning

\paragraph{Recorded answer/context.}
D: D: \$118\textbackslash{}

\paragraph{Figure/image path.}
/root/proposal\_for\_physic/science-mango/phyx\_mini\_run/phyx\_data/test\_image/240.png

\paragraph{Formalization target.}
create a compiling Lean file with sorry bodies at `PhyXMiniProblems/problem\_phyx\_mini\_0240.lean`. The Lean declarations must preserve the physical quantities, dimensions or dimensional roles, figure labels, governing-law hypotheses, and final relation expressed by this problem.
Use Mathlib/Physlib names found through LeanExplore where available. If a physics API is missing, introduce faithful local abstractions rather than scalar placeholder aliases.

\begin{theorem}[Physics formalization target]
\label{thm:physics:phyx_mini_0240:target}
\lean{PhyXMiniProblems.ProblemPhyXMini0240.maximumTransverseAcceleration}
\uses{def:physics:phyx-mini-0240:phyxminiproblems-problemphyxmini0240-amplitudeinmeters, def:physics:phyx-mini-0240:phyxminiproblems-problemphyxmini0240-angularfrequencyinradianspersecond, def:physics:phyx-mini-0240:phyxminiproblems-problemphyxmini0240-travelingstringwavesetup, def:physics:phyx-mini-0240:phyxminiproblems-problemphyxmini0240-matchesprimarystringwavefigure, def:physics:phyx-mini-0240:phyxminiproblems-problemphyxmini0240-hasstringwaveproblemdata, def:physics:phyx-mini-0240:phyxminiproblems-problemphyxmini0240-hasphysicalstringwaveparameters, def:physics:phyx-mini-0240:phyxminiproblems-problemphyxmini0240-satisfiestravelingstringwavelaws, def:physics:phyx-mini-0240:phyxminiproblems-problemphyxmini0240-transverseaccelerationmagnituderange, def:physics:phyx-mini-0240:phyxminiproblems-problemphyxmini0240-accelerationanswerchoice, def:physics:phyx-mini-0240:phyxminiproblems-problemphyxmini0240-displayedaccelerationinmeterspersecondsquared}
The assigned autoformalize agent should translate this physics problem into Lean declarations in the covered file, with theorem and lemma proof bodies written as `by sorry`.
\end{theorem}
\begin{proof}
This is an autoformalization task, not a proof task. Produce statements that can later be proved by the physics prover without weakening the physical model.
\end{proof}

% --- Archon declaration topology begin ---
\section{Declaration topology}
\begin{definition}[Amplitude Quantity]
  \label{def:physics:phyx-mini-0240:phyxminiproblems-problemphyxmini0240-amplitudequantity}
  \lean{PhyXMiniProblems.ProblemPhyXMini0240.AmplitudeQuantity}
  A nonnegative physical transverse amplitude
\end{definition}
\begin{proof}
  This is the declared model definition of Amplitude Quantity.
\end{proof}
\begin{definition}[Signed Length Quantity]
  \label{def:physics:phyx-mini-0240:phyxminiproblems-problemphyxmini0240-signedlengthquantity}
  \lean{PhyXMiniProblems.ProblemPhyXMini0240.SignedLengthQuantity}
  A signed physical axial position or transverse displacement
\end{definition}
\begin{proof}
  This is the declared model definition of Signed Length Quantity.
\end{proof}
\begin{definition}[Time Quantity]
  \label{def:physics:phyx-mini-0240:phyxminiproblems-problemphyxmini0240-timequantity}
  \lean{PhyXMiniProblems.ProblemPhyXMini0240.TimeQuantity}
  A signed physical time coordinate relative to a chosen origin
\end{definition}
\begin{proof}
  This is the declared model definition of Time Quantity.
\end{proof}
\begin{definition}[Cyclic Frequency Quantity]
  \label{def:physics:phyx-mini-0240:phyxminiproblems-problemphyxmini0240-cyclicfrequencyquantity}
  \lean{PhyXMiniProblems.ProblemPhyXMini0240.CyclicFrequencyQuantity}
  An ordinary cyclic frequency, whose SI readout is in hertz
\end{definition}
\begin{proof}
  This is the declared model definition of Cyclic Frequency Quantity.
\end{proof}
\begin{definition}[Angular Frequency Quantity]
  \label{def:physics:phyx-mini-0240:phyxminiproblems-problemphyxmini0240-angularfrequencyquantity}
  \lean{PhyXMiniProblems.ProblemPhyXMini0240.AngularFrequencyQuantity}
  An angular frequency, whose SI readout is in radians per second
\end{definition}
\begin{proof}
  This is the declared model definition of Angular Frequency Quantity.
\end{proof}
\begin{definition}[Wave Number Quantity]
  \label{def:physics:phyx-mini-0240:phyxminiproblems-problemphyxmini0240-wavenumberquantity}
  \lean{PhyXMiniProblems.ProblemPhyXMini0240.WaveNumberQuantity}
  A wave number, whose SI readout is in radians per meter
\end{definition}
\begin{proof}
  This is the declared model definition of Wave Number Quantity.
\end{proof}
\begin{definition}[Speed Quantity]
  \label{def:physics:phyx-mini-0240:phyxminiproblems-problemphyxmini0240-speedquantity}
  \lean{PhyXMiniProblems.ProblemPhyXMini0240.SpeedQuantity}
  A nonnegative physical propagation speed
\end{definition}
\begin{proof}
  This is the declared model definition of Speed Quantity.
\end{proof}
\begin{definition}[Acceleration Quantity]
  \label{def:physics:phyx-mini-0240:phyxminiproblems-problemphyxmini0240-accelerationquantity}
  \lean{PhyXMiniProblems.ProblemPhyXMini0240.AccelerationQuantity}
  A signed transverse acceleration component
\end{definition}
\begin{proof}
  This is the declared model definition of Acceleration Quantity.
\end{proof}
\begin{definition}[amplitude In Meters]
  \label{def:physics:phyx-mini-0240:phyxminiproblems-problemphyxmini0240-amplitudeinmeters}
  \lean{PhyXMiniProblems.ProblemPhyXMini0240.amplitudeInMeters}
  \uses{def:physics:phyx-mini-0240:phyxminiproblems-problemphyxmini0240-amplitudequantity}
  Read a nonnegative amplitude in meters
\end{definition}
\begin{proof}
  This is the declared model definition of amplitude In Meters.
\end{proof}
\begin{definition}[amplitude In Centimeters]
  \label{def:physics:phyx-mini-0240:phyxminiproblems-problemphyxmini0240-amplitudeincentimeters}
  \lean{PhyXMiniProblems.ProblemPhyXMini0240.amplitudeInCentimeters}
  \uses{def:physics:phyx-mini-0240:phyxminiproblems-problemphyxmini0240-amplitudequantity}
  Read a nonnegative amplitude in centimeters
\end{definition}
\begin{proof}
  This is the declared model definition of amplitude In Centimeters.
\end{proof}
\begin{definition}[signed Length In Meters]
  \label{def:physics:phyx-mini-0240:phyxminiproblems-problemphyxmini0240-signedlengthinmeters}
  \lean{PhyXMiniProblems.ProblemPhyXMini0240.signedLengthInMeters}
  \uses{def:physics:phyx-mini-0240:phyxminiproblems-problemphyxmini0240-signedlengthquantity}
  Read a signed axial position or displacement in meters
\end{definition}
\begin{proof}
  This is the declared model definition of signed Length In Meters.
\end{proof}
\begin{definition}[time In Seconds]
  \label{def:physics:phyx-mini-0240:phyxminiproblems-problemphyxmini0240-timeinseconds}
  \lean{PhyXMiniProblems.ProblemPhyXMini0240.timeInSeconds}
  \uses{def:physics:phyx-mini-0240:phyxminiproblems-problemphyxmini0240-timequantity}
  Read a time coordinate in seconds
\end{definition}
\begin{proof}
  This is the declared model definition of time In Seconds.
\end{proof}
\begin{definition}[frequency In Hertz]
  \label{def:physics:phyx-mini-0240:phyxminiproblems-problemphyxmini0240-frequencyinhertz}
  \lean{PhyXMiniProblems.ProblemPhyXMini0240.frequencyInHertz}
  \uses{def:physics:phyx-mini-0240:phyxminiproblems-problemphyxmini0240-cyclicfrequencyquantity}
  Read an ordinary cyclic frequency in hertz
\end{definition}
\begin{proof}
  This is the declared model definition of frequency In Hertz.
\end{proof}
\begin{definition}[angular Frequency In Radians Per Second]
  \label{def:physics:phyx-mini-0240:phyxminiproblems-problemphyxmini0240-angularfrequencyinradianspersecond}
  \lean{PhyXMiniProblems.ProblemPhyXMini0240.angularFrequencyInRadiansPerSecond}
  \uses{def:physics:phyx-mini-0240:phyxminiproblems-problemphyxmini0240-angularfrequencyquantity}
  Read an angular frequency in radians per second
\end{definition}
\begin{proof}
  This is the declared model definition of angular Frequency In Radians Per Second.
\end{proof}
\begin{definition}[wave Number In Radians Per Meter]
  \label{def:physics:phyx-mini-0240:phyxminiproblems-problemphyxmini0240-wavenumberinradianspermeter}
  \lean{PhyXMiniProblems.ProblemPhyXMini0240.waveNumberInRadiansPerMeter}
  \uses{def:physics:phyx-mini-0240:phyxminiproblems-problemphyxmini0240-wavenumberquantity}
  Read a wave number in radians per meter
\end{definition}
\begin{proof}
  This is the declared model definition of wave Number In Radians Per Meter.
\end{proof}
\begin{definition}[speed In Meters Per Second]
  \label{def:physics:phyx-mini-0240:phyxminiproblems-problemphyxmini0240-speedinmeterspersecond}
  \lean{PhyXMiniProblems.ProblemPhyXMini0240.speedInMetersPerSecond}
  \uses{def:physics:phyx-mini-0240:phyxminiproblems-problemphyxmini0240-speedquantity}
  Read a propagation speed in meters per second
\end{definition}
\begin{proof}
  This is the declared model definition of speed In Meters Per Second.
\end{proof}
\begin{definition}[acceleration In Meters Per Second Squared]
  \label{def:physics:phyx-mini-0240:phyxminiproblems-problemphyxmini0240-accelerationinmeterspersecondsquared}
  \lean{PhyXMiniProblems.ProblemPhyXMini0240.accelerationInMetersPerSecondSquared}
  \uses{def:physics:phyx-mini-0240:phyxminiproblems-problemphyxmini0240-accelerationquantity}
  Read a signed transverse acceleration in meters per second squared
\end{definition}
\begin{proof}
  This is the declared model definition of acceleration In Meters Per Second Squared.
\end{proof}
\begin{definition}[String Wave Kind]
  \label{def:physics:phyx-mini-0240:phyxminiproblems-problemphyxmini0240-stringwavekind}
  \lean{PhyXMiniProblems.ProblemPhyXMini0240.StringWaveKind}
  The disturbance carried by the string
\end{definition}
\begin{proof}
  This is the declared model definition of String Wave Kind.
\end{proof}
\begin{definition}[Propagation Direction]
  \label{def:physics:phyx-mini-0240:phyxminiproblems-problemphyxmini0240-propagationdirection}
  \lean{PhyXMiniProblems.ProblemPhyXMini0240.PropagationDirection}
  Direction of travel along the string's axial coordinate
\end{definition}
\begin{proof}
  This is the declared model definition of Propagation Direction.
\end{proof}
\begin{definition}[String Attachment]
  \label{def:physics:phyx-mini-0240:phyxminiproblems-problemphyxmini0240-stringattachment}
  \lean{PhyXMiniProblems.ProblemPhyXMini0240.StringAttachment}
  Mechanical attachment shown at the left end of the string
\end{definition}
\begin{proof}
  This is the declared model definition of String Attachment.
\end{proof}
\begin{definition}[String Wave Figure Label]
  \label{def:physics:phyx-mini-0240:phyxminiproblems-problemphyxmini0240-stringwavefigurelabel}
  \lean{PhyXMiniProblems.ProblemPhyXMini0240.StringWaveFigureLabel}
  Text labels visible in the primary figure
\end{definition}
\begin{proof}
  This is the declared model definition of String Wave Figure Label.
\end{proof}
\begin{definition}[String Wave Figure Feature]
  \label{def:physics:phyx-mini-0240:phyxminiproblems-problemphyxmini0240-stringwavefigurefeature}
  \lean{PhyXMiniProblems.ProblemPhyXMini0240.StringWaveFigureFeature}
  Geometric features to which the figure labels point
\end{definition}
\begin{proof}
  This is the declared model definition of String Wave Figure Feature.
\end{proof}
\begin{definition}[Traveling String Wave Setup]
  \label{def:physics:phyx-mini-0240:phyxminiproblems-problemphyxmini0240-travelingstringwavesetup}
  \lean{PhyXMiniProblems.ProblemPhyXMini0240.TravelingStringWaveSetup}
  \uses{def:physics:phyx-mini-0240:phyxminiproblems-problemphyxmini0240-amplitudequantity, def:physics:phyx-mini-0240:phyxminiproblems-problemphyxmini0240-signedlengthquantity, def:physics:phyx-mini-0240:phyxminiproblems-problemphyxmini0240-timequantity, def:physics:phyx-mini-0240:phyxminiproblems-problemphyxmini0240-cyclicfrequencyquantity, def:physics:phyx-mini-0240:phyxminiproblems-problemphyxmini0240-angularfrequencyquantity, def:physics:phyx-mini-0240:phyxminiproblems-problemphyxmini0240-wavenumberquantity, def:physics:phyx-mini-0240:phyxminiproblems-problemphyxmini0240-speedquantity, def:physics:phyx-mini-0240:phyxminiproblems-problemphyxmini0240-accelerationquantity, def:physics:phyx-mini-0240:phyxminiproblems-problemphyxmini0240-stringwavekind, def:physics:phyx-mini-0240:phyxminiproblems-problemphyxmini0240-propagationdirection, def:physics:phyx-mini-0240:phyxminiproblems-problemphyxmini0240-stringattachment, def:physics:phyx-mini-0240:phyxminiproblems-problemphyxmini0240-stringwavefigurelabel, def:physics:phyx-mini-0240:phyxminiproblems-problemphyxmini0240-stringwavefigurefeature}
  The independent physical quantities, figure roles, and time-dependent fields of the modeled string wave. The displacement and acceleration remain independent fields until the governing-law premise is supplied
\end{definition}
\begin{proof}
  This is the declared model definition of Traveling String Wave Setup.
\end{proof}
\begin{definition}[Matches Primary String Wave Figure]
  \label{def:physics:phyx-mini-0240:phyxminiproblems-problemphyxmini0240-matchesprimarystringwavefigure}
  \lean{PhyXMiniProblems.ProblemPhyXMini0240.MatchesPrimaryStringWaveFigure}
  \uses{def:physics:phyx-mini-0240:phyxminiproblems-problemphyxmini0240-travelingstringwavesetup}
  Qualitative information read from the primary figure. In particular, the A arrow measures crest height from equilibrium and the velocity arrow points right. This predicate contains no acceleration value
\end{definition}
\begin{proof}
  This is the declared model definition of Matches Primary String Wave Figure.
\end{proof}
\begin{definition}[Has String Wave Problem Data]
  \label{def:physics:phyx-mini-0240:phyxminiproblems-problemphyxmini0240-hasstringwaveproblemdata}
  \lean{PhyXMiniProblems.ProblemPhyXMini0240.HasStringWaveProblemData}
  \uses{def:physics:phyx-mini-0240:phyxminiproblems-problemphyxmini0240-amplitudeincentimeters, def:physics:phyx-mini-0240:phyxminiproblems-problemphyxmini0240-signedlengthinmeters, def:physics:phyx-mini-0240:phyxminiproblems-problemphyxmini0240-timeinseconds, def:physics:phyx-mini-0240:phyxminiproblems-problemphyxmini0240-frequencyinhertz, def:physics:phyx-mini-0240:phyxminiproblems-problemphyxmini0240-speedinmeterspersecond, def:physics:phyx-mini-0240:phyxminiproblems-problemphyxmini0240-travelingstringwavesetup}
  Numerical and origin data stated in the problem. The final field is the given condition y = 0 at x = 0 and t = 0, not a maximum-acceleration condition
\end{definition}
\begin{proof}
  This is the declared model definition of Has String Wave Problem Data.
\end{proof}
\begin{definition}[Has Physical String Wave Parameters]
  \label{def:physics:phyx-mini-0240:phyxminiproblems-problemphyxmini0240-hasphysicalstringwaveparameters}
  \lean{PhyXMiniProblems.ProblemPhyXMini0240.HasPhysicalStringWaveParameters}
  \uses{def:physics:phyx-mini-0240:phyxminiproblems-problemphyxmini0240-amplitudeinmeters, def:physics:phyx-mini-0240:phyxminiproblems-problemphyxmini0240-frequencyinhertz, def:physics:phyx-mini-0240:phyxminiproblems-problemphyxmini0240-angularfrequencyinradianspersecond, def:physics:phyx-mini-0240:phyxminiproblems-problemphyxmini0240-wavenumberinradianspermeter, def:physics:phyx-mini-0240:phyxminiproblems-problemphyxmini0240-speedinmeterspersecond, def:physics:phyx-mini-0240:phyxminiproblems-problemphyxmini0240-travelingstringwavesetup}
  Positivity and nondegeneracy of the physical wave parameters
\end{definition}
\begin{proof}
  This is the declared model definition of Has Physical String Wave Parameters.
\end{proof}
\begin{definition}[Satisfies Traveling String Wave Laws]
  \label{def:physics:phyx-mini-0240:phyxminiproblems-problemphyxmini0240-satisfiestravelingstringwavelaws}
  \lean{PhyXMiniProblems.ProblemPhyXMini0240.SatisfiesTravelingStringWaveLaws}
  \uses{def:physics:phyx-mini-0240:phyxminiproblems-problemphyxmini0240-signedlengthquantity, def:physics:phyx-mini-0240:phyxminiproblems-problemphyxmini0240-timequantity, def:physics:phyx-mini-0240:phyxminiproblems-problemphyxmini0240-amplitudeinmeters, def:physics:phyx-mini-0240:phyxminiproblems-problemphyxmini0240-signedlengthinmeters, def:physics:phyx-mini-0240:phyxminiproblems-problemphyxmini0240-timeinseconds, def:physics:phyx-mini-0240:phyxminiproblems-problemphyxmini0240-frequencyinhertz, def:physics:phyx-mini-0240:phyxminiproblems-problemphyxmini0240-angularfrequencyinradianspersecond, def:physics:phyx-mini-0240:phyxminiproblems-problemphyxmini0240-wavenumberinradianspermeter, def:physics:phyx-mini-0240:phyxminiproblems-problemphyxmini0240-speedinmeterspersecond, def:physics:phyx-mini-0240:phyxminiproblems-problemphyxmini0240-accelerationinmeterspersecondsquared, def:physics:phyx-mini-0240:phyxminiproblems-problemphyxmini0240-travelingstringwavesetup}
  Generic governing laws for a right-moving sinusoidal string wave: ω = 2πf; the nondispersive relation ω = v k; y(x,t) = A sin(k(x-x₀) - ω(t-t₀) + φ); each string element undergoes transverse SHM, a\_y = -ω² y. These pointwise laws do not state a maximum or select any answer choice
\end{definition}
\begin{proof}
  This is the declared model definition of Satisfies Traveling String Wave Laws.
\end{proof}
\begin{definition}[String Element State]
  \label{def:physics:phyx-mini-0240:phyxminiproblems-problemphyxmini0240-stringelementstate}
  \lean{PhyXMiniProblems.ProblemPhyXMini0240.StringElementState}
  \uses{def:physics:phyx-mini-0240:phyxminiproblems-problemphyxmini0240-signedlengthquantity, def:physics:phyx-mini-0240:phyxminiproblems-problemphyxmini0240-timequantity, def:physics:phyx-mini-0240:phyxminiproblems-problemphyxmini0240-signedlengthinmeters, def:physics:phyx-mini-0240:phyxminiproblems-problemphyxmini0240-travelingstringwavesetup}
  A state of an element on the modeled string. The axial coordinate is at or to the right of the wall attachment; time is unrestricted
\end{definition}
\begin{proof}
  This is the declared model definition of String Element State.
\end{proof}
\begin{definition}[transverse Acceleration Magnitude Range]
  \label{def:physics:phyx-mini-0240:phyxminiproblems-problemphyxmini0240-transverseaccelerationmagnituderange}
  \lean{PhyXMiniProblems.ProblemPhyXMini0240.transverseAccelerationMagnitudeRange}
  \uses{def:physics:phyx-mini-0240:phyxminiproblems-problemphyxmini0240-accelerationinmeterspersecondsquared, def:physics:phyx-mini-0240:phyxminiproblems-problemphyxmini0240-travelingstringwavesetup, def:physics:phyx-mini-0240:phyxminiproblems-problemphyxmini0240-stringelementstate}
  The set of all attained magnitudes of transverse acceleration, in SI
\end{definition}
\begin{proof}
  This is the declared model definition of transverse Acceleration Magnitude Range.
\end{proof}
\begin{definition}[Acceleration Answer Choice]
  \label{def:physics:phyx-mini-0240:phyxminiproblems-problemphyxmini0240-accelerationanswerchoice}
  \lean{PhyXMiniProblems.ProblemPhyXMini0240.AccelerationAnswerChoice}
  Labels of the four acceleration choices printed with the problem
\end{definition}
\begin{proof}
  This is the declared model definition of Acceleration Answer Choice.
\end{proof}
\begin{definition}[displayed Acceleration In Meters Per Second Squared]
  \label{def:physics:phyx-mini-0240:phyxminiproblems-problemphyxmini0240-displayedaccelerationinmeterspersecondsquared}
  \lean{PhyXMiniProblems.ProblemPhyXMini0240.displayedAccelerationInMetersPerSecondSquared}
  \uses{def:physics:phyx-mini-0240:phyxminiproblems-problemphyxmini0240-accelerationanswerchoice}
  The displayed answer values, interpreted as meter-per-second-squared readouts because the requested quantity is transverse acceleration
\end{definition}
\begin{proof}
  This is the declared model definition of displayed Acceleration In Meters Per Second Squared.
\end{proof}
\begin{definition}[recorded Dataset Answer]
  \label{def:physics:phyx-mini-0240:phyxminiproblems-problemphyxmini0240-recordeddatasetanswer}
  \lean{PhyXMiniProblems.ProblemPhyXMini0240.recordedDatasetAnswer}
  \uses{def:physics:phyx-mini-0240:phyxminiproblems-problemphyxmini0240-accelerationanswerchoice}
  The answer label recorded in the supplied dataset metadata
\end{definition}
\begin{proof}
  This is the declared model definition of recorded Dataset Answer.
\end{proof}
\begin{lemma}[angular Frequency In Radians Per Second eq ten pi]
  \label{lem:physics:phyx-mini-0240:phyxminiproblems-problemphyxmini0240-angularfrequencyinradianspersecond-eq-ten-pi}
  \lean{PhyXMiniProblems.ProblemPhyXMini0240.angularFrequencyInRadiansPerSecond_eq_ten_pi}
  \uses{def:physics:phyx-mini-0240:phyxminiproblems-problemphyxmini0240-angularfrequencyinradianspersecond, def:physics:phyx-mini-0240:phyxminiproblems-problemphyxmini0240-travelingstringwavesetup, def:physics:phyx-mini-0240:phyxminiproblems-problemphyxmini0240-hasstringwaveproblemdata, def:physics:phyx-mini-0240:phyxminiproblems-problemphyxmini0240-satisfiestravelingstringwavelaws}
  The stated frequency gives angular frequency ω = 10π rad/s
\end{lemma}
\begin{proof}
  The conclusion follows from the governing relations and figure readouts named in the statement of angular Frequency In Radians Per Second eq ten pi.
\end{proof}
\begin{lemma}[wave Number In Radians Per Meter eq pi over two]
  \label{lem:physics:phyx-mini-0240:phyxminiproblems-problemphyxmini0240-wavenumberinradianspermeter-eq-pi-over-two}
  \lean{PhyXMiniProblems.ProblemPhyXMini0240.waveNumberInRadiansPerMeter_eq_pi_over_two}
  \uses{def:physics:phyx-mini-0240:phyxminiproblems-problemphyxmini0240-wavenumberinradianspermeter, def:physics:phyx-mini-0240:phyxminiproblems-problemphyxmini0240-travelingstringwavesetup, def:physics:phyx-mini-0240:phyxminiproblems-problemphyxmini0240-hasstringwaveproblemdata, def:physics:phyx-mini-0240:phyxminiproblems-problemphyxmini0240-satisfiestravelingstringwavelaws}
  The supplied 20 m/s speed and 5 Hz drive imply the intermediate wave-number readout k = π/2 rad/m. This records the role of the given propagation speed even though the transverse-acceleration maximum is independent of it
\end{lemma}
\begin{proof}
  The conclusion follows from the governing relations and figure readouts named in the statement of wave Number In Radians Per Meter eq pi over two.
\end{proof}
% --- Archon declaration topology end ---
% --- Archon physics formalization source end ---
